\documentclass[12pt]{article}
\usepackage[utf8]{inputenc}
\usepackage[T1]{fontenc}
\usepackage{hyperref}
\title{Waling Language Resources for K-12 Curriculum Development}
\author{Compiled by Claude Sonnet 5.5 (claude-sonnet-5-5) for Andr\'as Kornai}
\date{}
\begin{document}
\maketitle

\section{Introduction}
Waling (ISO 639-3: wly; Glottocode wali1261) is a Kiranti (Bantawic) language of Sino-Tibetan, spoken in eastern Nepal. Glottolog marks it as dormant, with endangerment level 9, and ASJP describes it as extinct. See \url{https://glottolog.org/resource/languoid/iso/wly} and \url{https://asjp.clld.org/languages/WALING}. It should not be confused with Walung (Walungge), a Tibetic variety of Taplejung District, which is a different language.

\section{Primary Educational Resources}
None found.

\section{Linguistic and Reference Materials}
The searches found no grammar or dictionary of Waling. ASJP holds a word list, listed at \url{https://asjp.clld.org/languages/WALING}. Glottolog has the classification record.

\section{Religious and Community Resources}
None found.

\section{Digital and Audio Resources}
None found.

\section{Curriculum Development Context}
Almost no documentation was found, and there may be no speakers. A revitalization curriculum would depend on a short ASJP word list and on any archival work not located here. Working with related Bantawic languages and with the local community would be needed.
\end{document}